\section{Certified counting and the continuous operator}
\label{sec:certified}

Sections~\ref{sec:cert}--\ref{sec:subspace} certify quantities computed from
$\Hh$ in double precision, and Section~\ref{sec:fp} names two limitations: the
count $N(\sigma)$ is a combinatorial object read off a floating-point
factorization, and every statement refers to $\Hh$ rather than to
$\mathcal H$. This section attacks both with two constructions that turn out
to be cheap: a directed-rounding count, which removes the first limitation in
one dimension and in two whenever the potential is separable, and a
variational bound on the continuous eigenvalue, which removes half of the
second---the upper side---for every potential and needs no reference solution
and no interpolation constant.

\subsection{Directed rounding}

An interval $[a,b]$ represents every real in it; an operation returns
$\{x\circ y : x\in[a,b],\,y\in[c,d]\}$ rounded outward. The implementation is
a few dozen lines: each endpoint is moved one ulp outward with
\texttt{nextafter}, which is enough because IEEE\,754 arithmetic is correctly
rounded, so the exact result of a single operation lies within half an ulp of
the computed one. This is not higher-precision arithmetic---the precision is
still double---but every rounding is forced to the side that preserves the
statement.

What the resulting count certifies must be stated carefully, because it is
easy to overclaim. The certified matrix is the one the implementation actually
stores: diagonal entries $h^{-2}+V_i$, off-diagonals $-\tfrac12h^{-2}$, each
taken as an exact real number. A certified count is therefore a theorem about
\emph{that} symmetric tridiagonal matrix. Any difference between it and the
matrix a reader would write down from \eqref{eq:Hh} is a modelling question
(rounding of $h^{-2}$, of the potential samples), not a counting question.

\subsection{One dimension: pivots rather than determinants}

For a symmetric tridiagonal $T=\Hh-\sigma I$ with diagonal $a_k$ and
off-diagonal $b_k$, the number of eigenvalues below $\sigma$ is the number of
negative pivots in the $LDL^\top$ factorization, i.e.\ of the recurrence
\begin{equation}
  q_1 = a_1-\sigma,\qquad q_k = (a_k-\sigma) - \frac{b_{k-1}^2}{q_{k-1}} ,
  \label{eq:pivot}
\end{equation}
by Sylvester's law of inertia. The determinant recurrence
$d_k=(a_k-\sigma)d_{k-1}-b_{k-1}^2d_{k-2}$ computes the same sign pattern in
exact arithmetic, and we tried it first: it is unusable. In intervals the
relative width of $d_k$ exceeds $1$---so the interval contains zero and the
sign is undecidable---after $35$ to $45$ steps, on every one of the $363$
scan shifts at $N=2001$ (\texttt{runs/directed\_count/stats.json}). The
determinant cancels catastrophically because it alternates in sign and nearly
does so at the scale of the spectrum's bulk. Equation~\eqref{eq:pivot} avoids
this: when the pivot is resolved, the quotient $b^2/q_{k-1}$ is a small
correction to $a_k-\sigma$, and the interval stays narrow to the end of the
recurrence. The cost is one division per node, which must be refused when the
pivot interval straddles zero.

That refusal is the whole point. A count that cannot be decided is reported as
uncertified with the index of the offending pivot, and the caller must handle
it; it is never rounded to an integer. On three potentials at $N=2001$, over
$121$ shifts spanning the low spectrum, $361$ of $363$ shifts were certified
and every certified value equals both the floating-point count and the dense
eigen-decomposition (\texttt{runs/directed\_count/stats.json}). Both refusals
are at the scan's last shift, which in each scan coincides exactly with an
eigenvalue in double precision---precisely the case a rigorous implementation
must refuse---and the refusal names the offending pivot (index $999$). A
median count costs $5$\,ms, against
$0.4$\,ms for the floating-point version (\texttt{runs/\_bench/count\_timing.json})---a
factor of thirteen.

Two further uses follow. Bisecting on the certified count produces a
\emph{certified bracket} for a discrete eigenvalue; when the bisection
reaches the resolution of the recurrence it stops and returns the last
verified bracket instead of failing. This is what the driver
\texttt{scripts/continuous\_bound.py} reports in the $\lambda_h$ column of
Table~\ref{tab:continuous} for the one-dimensional rows and, through the
Kronecker route described below, for the two separable two-dimensional
potentials as well; only the anharmonic potential's
$\lambda_h$ is a sparse-factorization eigenvalue rather than a certified
bracket.

\subsection{Two dimensions: a Kronecker sum, and why not elimination}

The obvious two-dimensional analogue is a banded $LDL^\top$ in intervals. It
does not work, and the failure is quantitative rather than a matter of
implementation: the elimination divides by pivots and the interval width is
amplified at each step, so pivots straddle zero long before the factorization
ends. At $361$ unknowns, $6$ of $10$ shifts could not be counted; the relative
pivot width reached $0.93$ immediately before one of the refusals, and the
earliest refusal came at step $174$ of the elimination, where the width was
$0.44$ (\texttt{runs/directed\_count/stats.json}). Symmetric pivoting does
not repair this, because the difficulty is overestimation, not a zero pivot.

For a separable potential, however, no elimination is needed at all. If
$V(x,y)=V_x(x)+V_y(y)$ then
\begin{equation}
  \Hh = A_x\otimes I + I\otimes A_y
  \label{eq:kron}
\end{equation}
exactly, with $A_x,A_y$ the one-dimensional operators on the two grids, so
$\spec(\Hh)=\{\mu_i+\nu_j\}$ and
\begin{equation}
  N(\sigma) = \#\{(i,j): \mu_i+\nu_j < \sigma\}.
\end{equation}
Each one-dimensional eigenvalue is bracketed by the certified count of the
previous subsection, and the two-dimensional count is a comparison of sums. A pair whose sum interval straddles $\sigma$ makes the count refuse. That
is not a property of degenerate levels as such---any $\sigma$ inside the
certified bracket of a pair refuses---but in the cluster protocol the scan
places $\sigma$ on a level, and for the box, whose spectrum has
multiplicities, the midpoint of a degenerate pair is the eigenvalue
itself. On the harmonic and box potentials at $81\times81$, all
$25$ shifts of a scan were certified with zero disagreements against the
floating-point inertia count, at $18$\,ms per count. The non-separable
anharmonic potential is refused outright, with a reason, rather than
silently falling back to floating point. Bisecting on this count yields
certified brackets for the two-dimensional eigenvalues themselves---widths
of $1.4$--$6.0\times10^{-11}$ for the harmonic oscillator and
$3.7\times10^{-10}$--$2.5\times10^{-9}$ for the box between $41$ and $101$
points per side---which is where the two-dimensional $\lambda_h$ of
Table~\ref{tab:continuous} comes from.

\subsection{A variational bound for the continuous operator}
\label{sec:variational}

Every statement so far concerns $\Hh$. The gap to $\mathcal H$ was left as an
open item because the standard routes---an eigenvalue-perturbation estimate
with an explicit $O(h^2)$ constant, or a homotopy/Newton enclosure of the
elliptic problem---need regularity or pointwise bounds that a trained network
does not supply.

One direction needs none of that. By the min--max principle,
\begin{equation}
  \lambda_1(\mathcal H) = \min_{u\in H^1_0}\ \frac{\tfrac12\int u'^2 + \int V u^2}{\int u^2}
  \;\le\; R(u)
  \qquad\text{for every } u\in H^1_0 ,
  \label{eq:minmax}
\end{equation}
so \emph{any} conforming trial function yields an upper bound on the
continuous ground-state energy. The P1 (bilinear, in 2D) interpolant of a
sampled function is continuous and vanishes at the boundary, hence is in
$H^1_0$: the network's own eigenfunction, once interpolated, is a legitimate
trial function. No reference solution, no grid operator and no interpolation
constant enters. This is the precise sense in which the certificate now
reaches the continuous operator.

The arithmetic of \eqref{eq:minmax} must be distinguished from its
mathematics. In exact arithmetic the stiffness term is exact on our grids:
piecewise constant in one dimension, and a polynomial of degree two in each
variable in two dimensions, which the Gauss rule integrates exactly. The
potential term is integrated by Gauss quadrature on each element, which is
exact when $V u^2$ is a polynomial of low degree there---true for the
harmonic oscillator, the box, and, with the breakpoints supplied as element
boundaries, the finite well. Rather than trust the hypothesis, the code
evaluates the potential integral under two Gauss rules of different order and
refuses the bound whenever the two intervals fail to intersect; a kinked
potential whose kink lies inside an element is caught by that check. What the
code does \emph{not} do is evaluate the element integrals in interval
arithmetic: the accumulations are outward-rounded, with a padding that grows
with the number of terms, but the element-level arithmetic and the evaluation
of $V$ at the Gauss nodes are ordinary double precision. For the potentials
here that rounding is a few ulp, orders of magnitude below the padding the
summation already carries and eleven orders below the bracket it is compared
with; it is nevertheless a gap, and a user who needs it closed must supply an
interval evaluation of $V$. The one-sided statement---$\lambda_1(\mathcal H)
\le R(\psi_h)$---is then a theorem about the computed quotient up to that
rounding, and exactly a theorem about the quotient in exact arithmetic.

Table~\ref{tab:continuous} applies \eqref{eq:minmax} to the coarse and fine
ends of each of the six problems. Two columns matter. The bracket
$R(\psi_h)-\lambda_h$ contracts by the ratio that $O(h^2)$ predicts when the
grid is refined: $6.2$--$6.3$ for the three two-dimensional potentials between
$41$ and $101$ points per side (prediction $6.25$) and $56$--$64$ for the
harmonic oscillator and the infinite well between $501$ and $4001$ (prediction $64$;
the finite well has no such ratio, because its bracket changes sign, as
discussed below). At the
coarse grid the reference value lies between the two columns in all five
cases where it is known; at the fine grid in four of five, the exception
being the finite well, discussed below. For the
two-dimensional anharmonic potential, where no reference exists, the bracket
is a \emph{one-sided} statement: the true ground-state energy is at most the
tabulated $R(\psi_h)$.

\begin{table}[t]
\centering
\footnotesize
\setlength{\tabcolsep}{2pt}
\begin{tabular}{@{}llrrll@{}}
\toprule
Potential & $N$ & $\lambda_h$ & $R(\psi_h)$ & $\lambda^\star$ & in \\
\midrule
Harmonic & 501 & 0.4999819994 & 0.5000180013 & 0.5000000083 & \checkmark \\
Harmonic & 4001 & 0.4999997318 & 0.5000002813 & 0.5000000083 & \checkmark \\
Infinite well & 501 & 4.9347859593 & 4.9348184354 & 4.9348022005 & \checkmark \\
Infinite well & 4001 & 4.9348016381 & 4.9348024542 & 4.9348022005 & \checkmark \\
Finite well & 501 & -48.9981281228 & -48.9787820344 & -48.9810479595 & \checkmark \\
Finite well & 4001 & -48.9773754207 & -48.9809713144 & -48.9810479595 & $\times$ \\
\midrule
Harmonic (2D) & 41 & 0.9993927428 & 1.0021728615 & 1.0007820274 & \checkmark \\
Harmonic (2D) & 101 & 1.0005601465 & 1.0010042177 & 1.0007820274 & \checkmark \\
Box (2D) & 41 & 9.8645320539 & 9.8746788339 & 9.8696044011 & \checkmark \\
Box (2D) & 101 & 9.8687926854 & 9.8704161708 & 9.8696044011 & \checkmark \\
Anharmonic (2D) & 41 & 1.0962601793 & 1.1002100307 & -- & -- \\
Anharmonic (2D) & 101 & 1.0979179490 & 1.0985487577 & -- & -- \\
\bottomrule
\end{tabular}

\caption{Discrete eigenvalue $\lambda_h$ (the midpoint of a directed-rounding
bracket wherever the potential is separable---all one-dimensional rows, and
the harmonic and box in two dimensions; a sparse-factorization eigenvalue
for the anharmonic potential) and the
variational upper bound $R(\psi_h)$ of \eqref{eq:minmax} on the P1/Q1
interpolant of the discrete eigenvector. The final column asks whether the
reference value $\lambda^\star$ lies in $[\lambda_h,R(\psi_h)]$; ``--'' means
no reference is known. The finite well's failure at $N=4001$ is real---it
fails already at $N=2001$---and is discussed in the text.}
\label{tab:continuous}
\end{table}

\subsection{The two sides are not symmetric}

The upper bound is unconditional; the lower side is not, and one of our own
potentials shows why. The finite well has a discontinuous $V$, so the local
truncation error of the three-point stencil is $O(h)$---not $O(h^2)$---near the
jump, and its \emph{sign depends on where the jump falls relative to the
grid}. Across the twelve grids $N=1000,\dots,1011$ on the same well, the
discrete ground energy lies above the continuous value on six and below on
six, with errors up to $1.4\times10^{-2}$ and no monotone decay
(\texttt{runs/continuous\_bound/alignment.json}). Figure~\ref{fig:alignment}
shows the twelve signs.

\begin{figure}[t]
\centering
\begin{tikzpicture}
\begin{axis}[
  width=0.78\linewidth, height=0.40\linewidth,
  xlabel={grid points $N$}, ylabel={$\lambda_h-\lambda^\star$},
  ymin=-0.017, ymax=0.017,
  xtick={1000,1002,1004,1006,1008,1010},
  tick label style={font=\footnotesize},
  label style={font=\small},
  grid=major, grid style={dotted,gray!45},
  axis line style={gray!70},
]
\addplot+[ycomb, mark=*, mark size=1.1pt, thick, blue!55!black]
  table[col sep=comma, x=N, y=error]
  {figures/data/finite_well_alignment.csv};
\draw[black!55, dashed, line width=0.4pt]
  (axis cs:999.5,0) -- (axis cs:1011.5,0);
\end{axis}
\end{tikzpicture}
\caption{The finite well on twelve consecutive grids: the same potential,
only the grid moved. The discrete ground energy alternates above and below
the continuous value, with no trend in $N$, so the sign of the
discretization error is a property of the alignment of the jump with the
grid rather than of the grid size. Plotted from
\texttt{runs/continuous\_bound/alignment.json}; the $N=4001$ point discussed
in the text lies outside the window.}
\label{fig:alignment}
\end{figure}

The mechanism is visible on a single grid. At $N=4001$ the jump falls exactly
on a node, and the value that node receives---inside the well or outside---is
the whole error: keeping the grid and sampling that one point as $-50$ instead
of $0$ moves $\lambda_h$ by $7.4\times10^{-3}$, from $+3.67\times10^{-3}$ to
$-3.73\times10^{-3}$, straddling the true value. The convention used in the
table (the node outside, so the well is one grid point narrower) is what puts
$\lambda_h$ above $\lambda^\star$, and is why the finite-well bracket of
Table~\ref{tab:continuous} does not contain it at the finer grid. Neither
convention is more accurate; the discrete energy of a discontinuous potential
at an aligned grid is not a variationally consistent estimate of the
continuous one, and no amount of rounding control changes that.

The lesson is not that the count is wrong---it is certified, and it is a
certified statement about the stored matrix---but that a certified discrete
calculation does not by itself certify a continuous quantity when the method
is not variationally consistent. The same grid on a smooth potential behaves:
for the harmonic oscillator the discrete value is below the continuous one at
every grid we tried.

For the infinite well the lower side can also be made rigorous, and by an elementary
inequality rather than a numerical observation. On $[0,1]$ with $h=1/(n+1)$,
the eigenvalues of $\Hh=-\tfrac12\Delta_h$ are
\begin{equation}
  \lambda_j^{(h)} = \frac{2}{h^2}\sin^2\frac{j\pi h}{2},
  \qquad
  \lambda_j = \frac{j^2\pi^2}{2},
  \label{eq:sine}
\end{equation}
and $\sin\theta\le\theta$ gives $\lambda_j^{(h)}\le\lambda_j$ for every $j$ and
every $h$. In two dimensions the same argument applies along each axis, and
since the ordering of sums is preserved, $\lambda^{(h)}_j\le\lambda_j$ holds
for the entire spectrum of the two-dimensional box. There both ends are
provable: the sine inequality puts the discrete eigenvalue below
$\lambda^\star$, \eqref{eq:minmax} puts $\lambda^\star$ below $R(\psi_h)$, and
the certified bracket locates the discrete eigenvalue to
$2.5\times10^{-9}$---five orders of magnitude inside the margin
$8\times10^{-4}$ between the two bounds at the finer grid, so the tabulated
midpoint can be replaced by the bracket's lower endpoint without changing the
conclusion. For the other potentials
we do not have a lower bound, and we do not claim one: the honest summary is
that \eqref{eq:minmax} makes the continuous eigenvalue certifiable from above
for every potential---for the harmonic oscillator the discrete value is below
the continuous one at every grid we tried (\texttt{bound.json})---and from
below only where a comparison principle is available.

\subsection{What this changes}

The protocol of Section~\ref{sec:protocol} is unchanged, and so is the cluster
construction: what limits a certificate at a degenerate level is the
resolution at which two shifts can be separated, and wherever the count is
certified, directed rounding replaces ``double precision with guards'' by
``provable or refused'' at that same resolution. The experiments of Section~\ref{sec:results} still count in
floating point with the guards of Section~\ref{sec:counting}, at $0.4$\,ms
rather than $5$\,ms per count: the two agree on every shift at which the
certified count does not refuse, and the certified function takes the same
arguments when the guarantee rather than the speed is wanted. What does change
is the answer to the question a referee
should ask of any certified eigenvalue of a \emph{discretized} problem: how
far is this from the equation you actually wanted to solve? The upper bound
answers it with a number, for every trial function the practitioner already
has---including the trained one---and it is the only one of the two sides that
holds without a comparison principle.
